\chapter{Conditional Prediction and Generative Memory}
\label{app:conditional-memory}

A generative model can relate a cue to a different result. Evidence about the cue then changes the predicted result through that relation. This appendix realizes the general agent model of \cref{def:generative-model} with correlated recognition, derives its exact local inference, and separates that inference from a subsequent compression of the posterior. The construction begins with general probability laws. Gaussian mixture regression supplies an explicit finite-dimensional example \citep{Calinon2016}; its use here makes the model-presentation sector and the retained information explicit.

\section{A normalized cue and result model}

\constructionheading{Conditional prediction model}{cmem:model}
Fix a context. Let $\mathsf M,\mathsf X,\mathsf Y,\mathsf O$ be standard Borel spaces. Choose a proper presentation reference $r(dm)$, a normalized cue kernel $\mu_m(du)$, a normalized result kernel $K_m(u,dv)$, and a normalized record kernel $\mathcal L(u,do)$. All kernels are jointly measurable in their conditioning variables. Define
\begin{equation}
 P(dm,du,dv,do)=r(dm)\mu_m(du)K_m(u,dv)\mathcal L(u,do).
 \label{cmem:joint}
\end{equation}
Successive integration gives total mass one. In the agent notation, take $\mathsf K=\mathsf X\times\mathsf Y$, evaluator $G(m,d(u,v))=\mu_m(du)K_m(u,dv)$, and record kernel $\ell((u,v),m;do)=\mathcal L(u,do)$. Thus \cref{cmem:joint} is precisely $rG\ell$. It declares conditional independence of the result and record given the cue and presentation.

Assume $\mathcal L(u,do)=l(o\mid u)\lambda(do)$ for a fixed reference measure. At a specified record $o$, suppose
\begin{equation}
 \begin{aligned}
 0<Z_m(o)&:=\int l(o\mid u)\mu_m(du)<\infty
 &&r\text{-almost everywhere},\\
 0<Z(o)&:=\int Z_m(o)r(dm)<\infty.
 \end{aligned}
 \label{cmem:evidence}
\end{equation}
Bayes' rule gives the presentation posterior and component conditional
\begin{equation}
 w^o(dm)=\frac{Z_m(o)}{Z(o)}r(dm),\qquad
 \Pi_m^o(du,dv)=\frac{l(o\mid u)\mu_m(du)K_m(u,dv)}{Z_m(o)}.
 \label{cmem:posterior}
\end{equation}
The joint recognition law is $\zeta^o(dm,du,dv)=w^o(dm)\Pi_m^o(du,dv)$. It has the general disintegration of \cref{eq:agent-disintegration}. Dependence between presentation and contents is retained. Since $K_m(u,\mathsf Y)=1$, the presentation evidence depends on the cue and record mechanism, while the conditional result kernel determines what each selected presentation predicts.

The represented content can instead be the result alone. Disintegrate $\mu_m(du)K_m(u,dv)=G_Y(m,dv)R_m(v,du)$ and set
\begin{equation}
 \ell_Y(v,m;do)=\int\mathcal L(u,do)R_m(v,du).
 \label{cmem:result-model}
\end{equation}
Disintegration on standard Borel spaces supplies measurable versions, with arbitrary extensions on null sets \citep{Kallenberg2021}. The kernel $\ell_Y$ is normalized, and $rG_Y\ell_Y$ is the $(M,Y,O)$ marginal of \cref{cmem:joint}. Both representations therefore give the same result posterior. A presentation-dependent likelihood is allowed by \cref{def:generative-model}.

\propositionheading{Correlated local free energy}{cmem:free-energy}
At a record satisfying \cref{cmem:evidence}, let a candidate joint recognition law be $\mathcal Q(dm,dk)=s(dm)q_m(dk)$, where $k=(u,v)$. Assume the absolute continuity, finite relative entropies, and integrable log likelihood needed below. Then
\begin{equation}
 \mathcal F(\mathcal Q;o)
 =\KL(s\Vert w^o)+\int\KL(q_m\Vert\Pi_m^o)s(dm)-\log Z(o).
 \label{cmem:free-energy-split}
\end{equation}
Consequently the unrestricted minimum is $-\log Z(o)$, attained uniquely by the posterior up to null sets.
\begin{proof}
The unnormalized joint at the observed record is $Z(o)\zeta^o$. Substituting this into the logarithmic density ratio defining free energy gives $\mathcal F=\KL(\mathcal Q\Vert\zeta^o)-\log Z(o)$. Disintegration splits the joint relative entropy into the two terms of \cref{cmem:free-energy-split}; their nonnegativity and equality conditions prove the minimum statement.
\end{proof}

For a finite presentation space with $r(a)=\hat\pi_a>0$, fixing the content conditionals gives
\begin{equation}
 s_a\propto\hat\pi_a e^{-f_a},\qquad
 f_a=\KL(q_a\Vert G_a)-\E_{q_a}\log l(o\mid X).
 \label{cmem:gibbs}
\end{equation}
At the exact conditionals, $f_a=-\log Z_a(o)$. These are Gibbs weights for optimized component free energies. They are distinct from the peer-source rows of \cref{ch:dynamics}: the label here chooses a model presentation.

The product restriction changes the inference problem. For the structured family $s_aq_X(du)K_a(u,dv)$, the result conditional cancels from the density ratio, and the label update is
\begin{equation}
 s_a\propto\hat\pi_a\exp[-\KL(q_X\Vert\mu_a)].
 \label{cmem:product-label}
\end{equation}
The record term is common to the labels at fixed $q_X$. With cue densities $\mu_a(u)$, varying $q_X$ gives $q_X(u)\propto l(o\mid u)\exp(\sum_a s_a\log\mu_a(u))$, whenever the normalizer is finite. This is a product over $(A,X)$, while retaining the result conditional; it is not a product over $(A,(X,Y))$. It illustrates the expected-log restriction described in \cref{rem:dyn-product} and classical mean-field inference \citep{BleiEtAl2017}. Exact conditioning, this structured approximation, and the product-family Fisher flow are separately specified operations.

\section{A Gaussian realization and its response}

Let $\mathsf M=\{1,\ldots,M\}$ and $\mathsf X=\mathsf Y=\mathsf O=\R^d$. Fix $\hat\pi_a>0$, $\sum_a\hat\pi_a=1$, vectors $\nu_a,\omega_a$, positive-definite matrices $L_a,Q_a,C$, and arbitrary matrices $F_a$. Set
\begin{equation}
 \begin{aligned}
 X\mid A=a&\sim\mathcal N(\nu_a,L_a),\\
 Y\mid X=u,A=a&\sim\mathcal N(\omega_a+F_a(u-\nu_a),Q_a),\\
 O\mid X=u&\sim\mathcal N(u,C).
 \end{aligned}
 \label{cmem:gaussian-model}
\end{equation}
The record covariance is fixed side information for this model. Each evaluator is the joint Gaussian with mean $(\nu_a,\omega_a)$ and covariance
\begin{equation}
 \begin{pmatrix}
 L_a&L_aF_a^T\\ F_aL_a&F_aL_aF_a^T+Q_a
 \end{pmatrix}.
 \label{cmem:joint-covariance}
\end{equation}
Its Schur complement is $Q_a\succ0$, so arbitrary, including singular, regressions are allowed.

\propositionheading{Exact component and result posteriors}{cmem:gaussian-posterior}
At $O=\xi$, put $W_a=L_a+C$. Then
\begin{equation}
 \begin{aligned}
 Z_a(\xi)&=\mathcal N(\xi;\nu_a,W_a),&
 w_a&=\frac{\hat\pi_a Z_a(\xi)}{\sum_b\hat\pi_bZ_b(\xi)},\\
 m_{X,a}&=\nu_a+L_aW_a^{-1}(\xi-\nu_a),&
 P_{X,a}&=(L_a^{-1}+C^{-1})^{-1},\\
 m_{Y,a}&=\omega_a+F_a(m_{X,a}-\nu_a),&
 P_{Y,a}&=Q_a+F_aP_{X,a}F_a^T.
 \end{aligned}
 \label{cmem:gaussian-posterior-formula}
\end{equation}
The result marginal is $p_Y^\xi=\sum_aw_a\mathcal N(m_{Y,a},P_{Y,a})$, with moments
\begin{equation}
 \bar m=\sum_aw_am_{Y,a},\qquad
 S=\sum_aw_a\left[P_{Y,a}+(m_{Y,a}-\bar m)(m_{Y,a}-\bar m)^T\right].
 \label{cmem:result-moments}
\end{equation}
\begin{proof}
The cue and independent record noise sum to a Gaussian with covariance $W_a$, giving $Z_a$. Completing the cue's quadratic log density gives precision $L_a^{-1}+C^{-1}$ and information vector $L_a^{-1}\nu_a+C^{-1}\xi$. Solving for the mean gives the displayed gain. Applying the affine result kernel gives its conditional moments and cross-covariance $\operatorname{Cov}(X,Y\mid\xi,a)=P_{X,a}F_a^T$. Total expectation and total covariance give \cref{cmem:result-moments}. These are conditional-mixture identities, as in Gaussian mixture regression \citep{Calinon2016}.
\end{proof}

For the result-only representation \cref{cmem:result-model}, let $V_a=F_aL_aF_a^T+Q_a$, $B_a=L_aF_a^TV_a^{-1}$, and $S_a=L_a-L_aF_a^TV_a^{-1}F_aL_a$. Then
\begin{equation}
 \begin{aligned}
 G_Y(a,dv)&=\mathcal N(dv;\omega_a,V_a),\\
 \ell_Y(v,a;do)&=\mathcal N(do;\nu_a+B_a(v-\omega_a),C+S_a).
 \end{aligned}
 \label{cmem:result-only-gaussian}
\end{equation}
Conditioning the joint evaluator on $Y=v$ gives cue covariance $S_a\succ0$, and adding independent record noise proves the likelihood formula.

\propositionheading{Cross-prediction response and posterior Fisher tensor}{cmem:response}
Hold the prior joint law of $(A,X,Y)$ and $C\succ0$ fixed, with record $O=X+\epsilon$ for independent $\epsilon\sim\mathcal N(0,C)$. Suppose the posterior moments used below exist and differentiation can be passed through their integrals on an open set of records. Then
\begin{equation}
 D_\xi\E[Y\mid\xi]=\operatorname{Cov}(Y,X\mid\xi)C^{-1},\qquad
 I_\xi=C^{-1}\operatorname{Cov}(X\mid\xi)C^{-1},
 \label{cmem:response-fisher}
\end{equation}
where $I_\xi$ is the Fisher tensor of the complete posterior family in record coordinates. These assumptions hold for the finite Gaussian model above.
\begin{proof}
The gradient of the log record likelihood is $C^{-1}(X-\xi)$. Subtracting its posterior expectation, the gradient of the log normalizer, gives posterior score $C^{-1}(X-\E[X\mid\xi])$. Multiplying the score by $Y$ and integrating proves the response identity. Its score second moment gives $I_\xi$. Finite Gaussian mixtures have smooth finite component moments and locally dominated Gaussian integrands, which justify these operations.
\end{proof}

For same-state prediction $Y=X$, the whitened Jacobian is $C^{-1/2}\operatorname{Cov}(X\mid\xi)C^{-1/2}$, positive definite in the Gaussian example. For a distinct result the cross-covariance can have either sign and need not be symmetric. Already with one component, $D_\xi\E[Y\mid\xi]=FL(L+C)^{-1}$. For any prescribed dense response $H$, choosing $F=H(L+C)L^{-1}$ gives $H$ exactly. The result center sets the affine intercept. This conditional response is not a gradient step on the cue, and its sign places no conflict with the positive Fisher tensor. The identity is a partial derivative at fixed bank and record covariance; it does not differentiate a larger computation in which those data also vary.

\section{Retained moments and statistical geometry}

\propositionheading{Gaussian projection and the loss of dependence}{cmem:projection}
For the positive finite Gaussian mixture $p_Y^\xi$ above, $q_G=\mathcal N(\bar m,S)$ uniquely minimizes $\KL(p_Y^\xi\Vert q)$ over nondegenerate Gaussians. If $\Pi_{A,Y}^\xi$ is the exact label/result joint, then
\begin{equation}
 \KL(\Pi_{A,Y}^\xi\Vert w\otimes q_G)
 =I_{\Pi^\xi}(A;Y)+\KL(p_Y^\xi\Vert q_G).
 \label{cmem:projection-loss}
\end{equation}
For the diagonal carrier $q_D=\mathcal N(\bar m,\operatorname{diag}S)$,
\begin{equation}
 \KL(p_Y^\xi\Vert q_D)-\KL(p_Y^\xi\Vert q_G)
 =\frac12\log\frac{\prod_k S_{kk}}{\det S}.
 \label{cmem:diagonal-loss}
\end{equation}
\begin{proof}
The expectation of a Gaussian log density depends only on the mean and covariance of the sampling law. Consequently, for every nondegenerate Gaussian $q$,
$\KL(p_Y^\xi\Vert q)-\KL(p_Y^\xi\Vert q_G)=\KL(q_G\Vert q)$.
This proves the minimum and its uniqueness. Inserting $w\otimes p_Y^\xi$ into the joint logarithmic ratio proves \cref{cmem:projection-loss}. Subtracting the two Gaussian log densities and using $\E[(Y-\bar m)(Y-\bar m)^T]=S$ gives \cref{cmem:diagonal-loss}. All entropies and relative entropies here are finite for the stated positive finite mixtures.
\end{proof}

This inclusive-divergence projection is the local moment-matching operation used in assumed-density filtering \citep{Minka2001EP}. It generally differs from minimizing the reverse divergence of variational free energy. For example, let $p_a=\tfrac12\mathcal N(-a,1)+\tfrac12\mathcal N(a,1)$ and $q_a=\mathcal N(0,1+a^2)$. The component $\mathcal N(a,1)$ has reverse divergence at most $\log2$ from $p_a$, whereas
\begin{equation}
 \KL(q_a\Vert p_a)
 \ge a^2-a\sqrt{\frac{2(1+a^2)}{\pi}}-\frac12\log(1+a^2)
 \longrightarrow\infty.
 \label{cmem:reverse-counterexample}
\end{equation}
Indeed, $p_a(y)\le(2\pi)^{-1/2}e^{-(|y|-a)^2/2}$; integrate its negative logarithm under $q_a$ and subtract the Gaussian entropy. Thus moment matching does not even minimize the reverse divergence for sufficiently separated components.

\propositionheading{Moments need not preserve later Gaussian records}{cmem:readout-loss}
For $0<t<1$, the admissible result mixture
\begin{equation}
 p_t=\tfrac12\mathcal N(t,1-t^2)+\tfrac12\mathcal N(-t,1-t^2)
 \label{cmem:constant-moments}
\end{equation}
has constant Gaussian carrier $q_G=\mathcal N(0,1)$, but a later record kernel $K_c(y,do)=\mathcal N(do;y,c)$, $c>0$, distinguishes $p_t$ from $q_G$.
\begin{proof}
Identical cue components, equal weights, $F_1=F_2=0$, result centers $\pm t$, and result noise $1-t^2$ realize $p_t$. Its mean is zero and variance one. Convolution with the record noise gives densities at zero
\begin{equation}
 (p_tK_c)(0)=\frac{e^{-t^2/[2(1-t^2+c)]}}{\sqrt{2\pi(1-t^2+c)}},\qquad
 (q_GK_c)(0)=\frac1{\sqrt{2\pi(1+c)}}.
\end{equation}
They agree at $t=0$, while the derivative of the logarithm of the first is $-t^3/(1-t^2+c)^2<0$. Their densities therefore differ for every $t>0$; continuity gives a differing probability on a small interval about zero.
\end{proof}

Moments preserve expectations of affine and quadratic readouts. \Cref{cmem:readout-loss} shows that this does not imply preservation of every later observation law. It concerns the chosen projection, not an obstruction to richer retained states or to general-law parents. Exact inheritance still requires the complete configuration retention and interface of \cref{inh:two-retentions,inh:deterministic-descent}.

For a regular family $\zeta_\theta(a,dv)=w_a\mathcal N(dv;m_a,P_a)$, the score split underlying \cref{geo:joint-fisher} gives
\begin{equation}
 \begin{split}
 g_\zeta^F(u,v)={}&\sum_a\frac{(d_uw_a)(d_vw_a)}{w_a}\\
 &+\sum_aw_a\left[(d_um_a)^TP_a^{-1}(d_vm_a)
 +\frac12\Tr(P_a^{-1}d_uP_aP_a^{-1}d_vP_a)\right].
 \end{split}
 \label{cmem:joint-fisher}
\end{equation}
The mixed term vanishes because each conditional Gaussian score has conditional mean zero. Null parameter directions must be removed before calling this tensor a metric. For the label/result family producing \cref{cmem:constant-moments}, its coefficient along $t$ is $(1+t^2)/(1-t^2)^2>0$, while the Gaussian carrier has zero coefficient. The result marginal itself also varies, with fourth moment $3-2t^4$. Joint statistical geometry therefore cannot be recovered from the moment carrier in general. These are admissible model families, not assertions about a particular evolution.

\section{Frames, sharing, and sequential records}

One group $\GL(d)$ acts on the doubled content by $h\cdot(u,v)=(hu,hv)$ and trivially on the finite labels. Re-expressing the evaluator and record data requires
\begin{equation}
 \begin{gathered}
 \nu_a'=h\nu_a,\quad \omega_a'=h\omega_a,\quad
 L_a'=hL_ah^T,\quad Q_a'=hQ_ah^T,\\
 F_a'=hF_ah^{-1},\quad C'=hCh^T,\quad \xi'=h\xi.
 \end{gathered}
 \label{cmem:frame-laws}
\end{equation}
The evidence densities acquire the same factor $|\det h|^{-1}$, so the weights are invariant; the conditional means and covariances transform by $h$ and congruence. This follows directly by substitution into \cref{cmem:gaussian-posterior-formula}. In local token frames $U_i$, the reference cue $(U_i^{-1}\mu_i,U_i^{-1}\Sigma_iU_i^{-T})$ is invariant under a simultaneous passive change of frame and local moments. A singular $F_a$ still defines a normalized response kernel; it is not a group-valued dictionary.

Sharing a bank supplies common reference and evaluator data. Sharing a generative model in the sense of \cref{def:shared-model} also requires matched record kernels under the declared frame maps. Different covariances $C_i$ instead give a shared family of conditional mechanisms indexed by supplied side information. The distinction must be stated if uncertainty varies between agents. A diagonal reference parametrization is not closed under arbitrary rebasing: if $hDh^T$ is diagonal for every positive diagonal $D$, its off-diagonal entries force $h_{ik}h_{jk}=0$ for $i\ne j$, so invertibility makes $h$ a permutation times an invertible diagonal matrix. This restriction on the reference parametrization does not remove local passive covariance with that reference fixed.

\constructionheading{An exact carried-presentation filter}{cmem:sequential}
Let the finite label evolve by a fixed row-stochastic matrix $T$. At step $l$, draw a fresh cue from $\mu_{A_l}$, a result from $K_{A_l}$, and a record from $\mathcal L_{C_l}$, with draws at distinct steps conditionally independent given the label sequence. For fixed positive record covariances, the finite law is
\begin{equation}
 \begin{split}
 &P(a_{0:n},du_{0:n},dv_{0:n},do_{0:n})\\
 &\quad=\hat\pi_{a_0}\prod_{l=1}^nT_{a_{l-1}a_l}
 \prod_{l=0}^n\left[\mu_{a_l}(du_l)K_{a_l}(u_l,dv_l)
 \mathcal L_{C_l}(u_l,do_l)\right].
 \end{split}
 \label{cmem:history}
\end{equation}
Integration and summation of the last step give one conditional on the previous label, proving normalization and consistency. If $w_l(a)=P(A_l=a\mid O_{0:l})$, marginalization and Bayes' rule give
\begin{equation}
 v_l(b)=\sum_aw_{l-1}(a)T_{ab},\qquad
 w_l(b)=\frac{v_l(b)Z_b(o_l;C_l)}{\sum_c v_l(c)Z_c(o_l;C_l)}.
 \label{cmem:filter}
\end{equation}
This is a hidden Markov emission model. The label vector suffices for its next label prediction because the next cue is fresh given the next label. A later interaction that reads the current result or its dependence on the label requires a different sufficiency argument. The continuous result law can still be needed for current readouts even in this model.

A modified prior proportional to $\hat\pi_bv_l(b)^\vartheta$ has a separate one-step interpretation. With baseline prior $\hat\pi$ and a record $o\mid B=b\sim\operatorname{Dir}(\alpha\mathbf1+\vartheta e_b)$, where $\alpha>0$ and $\vartheta\ge0$, its normalizer is independent of $b$, and the density at an interior simplex value $o=v_l$ is proportional to $v_l(b)^\vartheta$. Conditioning gives the modified prior. It equals the prediction in \cref{cmem:filter} for uniform $\hat\pi$ and $\vartheta=1$. Combining this record with a cue requires declaring their joint conditional law; the algebra alone does not supply conditional independence.

\Cref{app:records} distinguishes normalization of a history from exact Bayesian assimilation under its record mechanism. That distinction applies here. If $O_2=f(O_1)$ and $O_1$ is already registered, conditioning on $(O_1,O_2)$ is the same as conditioning on $O_1$; an independent-replicate likelihood would count the evidence again. Communication may bring information that is new to one receiver even in a deterministic population computation. A repeated conditional-prediction agent must therefore specify its accessible history, communication-generated record, retained joint law, and any projection. The local identities above establish neither a common free energy for an arbitrary composition nor a new formation result.
